\chapter{Project Overview and Assignment}
\label{ch:overview}

\section{The engineering problem}
The EN3563 mini-project asks a group to implement a pick-and-place task with a
robot arm that has at least four degrees of freedom. The work is intentionally
end-to-end: model the robot, assign link coordinate frames, construct a DH
table, compute forward kinematics (FK), formulate inverse kinematics (IK),
plan a smooth trajectory, control the joints, keep the carried object upright,
and submit a report and screen-recorded demonstration. The supplied sheet
offers physical construction or a Gazebo simulation using the supplied
Yaskawa GP7 CAD. The simulation option explicitly allows a simplified gripper
because the supplied CAD robot does not include an end effector. It specifies
the robot at the origin, a large Box A of $0.30\times0.30\times0.20$ m, a
$0.05$ m cube Box B at a random reachable ground location, and Box B placement
at the center of Box A's top surface while preserving orientation
\cite{assignment}.

The implemented system is a ROS 2 and Gazebo Sim application around a six-axis
GP7 URDF. It is more than a static kinematics exercise: the launch code creates
a task-specific robot description, starts Gazebo and the controllers, chooses
Box B's pickup pose, and starts a Python node that plans, executes, monitors,
attaches, releases, and verifies the object.

\section{Original requirement, clarification, and implementation}
The original sheet places Box A at $(1,1,0)$ m. Relative to the robot base at
the origin, this has horizontal radius
\begin{equation}
  r_A=\sqrt{1^2+1^2}=\sqrt{2}\approx\SI{1.414}{m}.
  \label{eq:box-a-radius}
\end{equation}
The repository's engineering notes estimate an upper horizontal reach below
\SI{1.03}{m} for this robot model, before joint-limit constraints. Thus the
printed destination is not achievable in the implemented scene. After the
sheet was issued, the lecturer verbally clarified that a reachable Box A
location could be used. The implementation consequently fixes Box A's ground
position at $(0.55,0.35,0)$ m. This is a later verbal clarification; it is not
an amendment printed in the supplied sheet, and the group has no written
permission to claim.

The assignment's random reachable ground pose refers to the movable small
object, Box B. The project keeps Box A fixed while offering two Box B pickup
modes:
\begin{itemize}
  \item \textbf{Far mode}: numerical constrained optimization finds a
  maximum-radius, downward-facing grasp. The actual pickup is moved inward by
  a \SI{5}{mm} margin. It is optimized and deterministic, not random.
  \item \textbf{Near mode}: the launch file samples $x\sim U(0.50,0.60)$ and
  $y\sim U(-0.23,-0.15)$ using Python's seeded pseudorandom generator. With
  seed 7 the observed pose was $(0.532383276,-0.217932066,0.025)$ m. This
  changes Box B's pose only. A sampled point is not certified reachable solely
  by the sampling bounds; the path's IK preflight can reject an infeasible
  sample.
\end{itemize}
The fixed destination and the two pickup definitions are plotted in
\cref{fig:scene}.

\begin{figure}[htbp]
  \centering
  \begin{tikzpicture}[>=Stealth,scale=4.1]
  \draw[->,gray,thick] (-.30,0)--(.82,0) node[right]{world $+x$ (m)};
  \draw[->,gray,thick] (0,-1.05)--(0,.58) node[above]{world $+y$ (m)};
  \fill[blue!12] (.40,.20) rectangle (.70,.50);
  \draw[blue!70!black,thick] (.40,.20) rectangle (.70,.50);
  \node[blue!50!black,align=center] at (.55,.35){fixed Box A\\$(.55,.35)$};
  \fill[red!75!black] (-.06874,-.92808) circle (.012);
  \node[red!65!black,anchor=east,align=right] at (-.10,-.928){far-mode B\\$(-.06874,-.92808)$};
  \fill[orange!85!black] (.53238,-.21793) circle (.012);
  \node[orange!65!black,anchor=west,align=left] at (.55,-.218){near mode, seed 7\\$(.53238,-.21793)$};
  \fill[black] (0,0) circle (.009) node[above left]{base $(0,0)$};
  \draw[->,black,dashed] (-.06874,-.92808)--(.55,.35);
  \draw[->,black,dashed] (.53238,-.21793)--(.55,.35);
\end{tikzpicture}

  \caption{World-plane schematic of the fixed Box A and example Box B pickup
  locations. The vertical coordinate is suppressed; both pickup centers are
  at $z=0.025$ m.}
  \label{fig:scene}
\end{figure}

Box A's center is at $z=0.10$ m because its height is \SI{0.20}{m}. The target
center for Box B is
\begin{equation}
 \pvec_{B,d}=\begin{bmatrix}0.55&0.35&0.225\end{bmatrix}^{T}\si{m},
 \label{eq:place-target}
\end{equation}
which places the cube's lower face at $z=0.20$ m, the top of Box A. In the
actual path the tool lowers to $z=0.226$ m, releases about \SI{1}{mm} above
the target center, and the cube settles under gravity before final pose
measurement.

\section{Requirements and evidence status}
\begin{longtable}{p{.29\textwidth}p{.31\textwidth}p{.32\textwidth}}
\caption{Assignment requirement mapped to implementation and evidence.}\label{tab:req-map}\\
\toprule Requirement & Implementation & Status / limitation\\\midrule\endfirsthead
\toprule Requirement & Implementation & Status / limitation\\\midrule\endhead
At least four degrees of freedom & Six revolute GP7 joints & Implemented in URDF; all six are commanded.\\
Coordinate frame diagram & Standard-DH frames & Derived and independently compared to URDF FK; see \cref{ch:dh}.\\
DH table and FK & Standard DH table plus tool transform & 104-config comparison recorded; runtime FK itself parses URDF, not DH constants.\\
IK for target pose & Bounded SciPy least squares & Numerical solver with URDF limits; no analytic GP7 solver.\\
Smooth trajectory / velocity control & Quintic blend, sampled Cartesian IK, trajectory action & Joint-speed cap is sampled; no measured acceleration or jerk report.\\
Random reachable Box B & Seeded near-mode coordinate sampling & Seed 7 was tested; arbitrary samples are preflighted and can fail.\\
Box A position & Fixed $(0.55,0.35,0)$ m & Used under lecturer's subsequent verbal clarification; original sheet says $(1,1,0)$.\\
Upright object transfer & Fixed downward grasp rotation; orientation monitored from Gazebo pose samples & Successful simulation records; sampling is not a continuous-time proof.\\
Gripper / grasp & Adaptive mesh gripper and detachable fixed-joint plugin & Simulated close and attachment; no demonstrated frictional or physical grasp.\\
Report and video & This guide and prior report draft; recording checklist & Group member fields and a new screen recording remain submission work.\\
\bottomrule
\end{longtable}

\section{What is complete versus what remains}
The repo's validation notes record successful builds of all three ROS packages,
six offline regression tests, 104 DH-versus-URDF FK configurations, and
headless Gazebo far and seeded-near runs. The executed sequence included
approach, descent, close/attach, lift, transfer, placement, release, and
retreat. Those are simulation results with an idealized fixed-joint retention
mechanism. No physical robot run, force-based grasp, CAD equivalence study, or
submission video is in the repository evidence. Chapter 13 preserves the
recorded metrics and separates them from any verification done while preparing
this guide.

\section{Scope boundaries}
This guide explains both mathematical background and the actual choices made
in source. It does not claim the project uses MoveIt, an analytic IK solution,
a collision planner, physical sensors, force control, or calibrated link
inertias. None of those appears in the working task path. ROS 2, Gazebo, and
the Python stack interact, but the repository has no custom C++ source file;
the C++ plugins and controller implementations are supplied by installed ROS
and Gazebo packages.
